\section{Lingkungan Implementasi}
\label{sec:lingkungan-implementasi}

Implementasi sistem dilakukan pada lingkungan prototipe yang merepresentasikan komponen utama DecMed, yaitu \textit{patient client}, \textit{hospital client}, \textit{ministry client}, \textit{Server} PRE, penyimpanan \textit{off-chain} IPFS, Redis, IOTA \textit{Gas Station}, dan \textit{smart contract} IOTA Move pada jaringan IOTA Localnet. Lingkungan ini dipilih agar rancangan pada Bab \ref{chap:analisis-dan-rancangan} dapat diuji tanpa bergantung pada integrasi langsung dengan sistem informasi rumah sakit yang sesungguhnya.

Pada implementasi DecMed sebelumnya, semua komponen dijalankan pada lingkungan lokal. Semua \textit{client}, \textit{Server} PRE, Redis PRE, kluster IPFS, IOTA \textit{Gas Station}, dan jaringan IOTA jenis Localnet dijalankan langsung di komputer lokal. Perangkat lokal yang digunakan tidak cukup kuat untuk menjalankan semua komponen secara bersamaan. Oleh karena itu, dilakukan pembagian komponen implementasi tersebut ke dalam lingkungan lokal dan dua buah \textit{Virtual Private Server} (VPS).

Komponen \textit{client} dan \textit{Server} PRE dijalankan pada lingkungan lokal. Lingkungan lokal digunakan untuk memudahkan pengembangan fitur, pengujian alur akses, serta pengamatan respons sistem. Komponen pendukung yang membutuhkan ketersediaan jaringan, yaitu kluster IPFS dan IOTA \textit{Gas Station}, dijalankan pada VPS. \textit{Smart contract} IOTA Move di-\textit{deploy} pada jaringan IOTA Localnet yang juga berjalan di VPS. IOTA Localnet digunakan karena sudah memadai untuk menyimulasikan interaksi dengan jaringan IOTA. VPS digunakan untuk menjalankan komponen pendukung karena komponen tersebut tidak banyak berubah selama implementasi, tetapi tetap dibutuhkan selama pengembangan. Pemanfaatan VPS juga mengurangi beban perangkat lokal karena perangkat lokal tidak perlu menjalankan seluruh komponen pendukung.

Berikut adalah spesifikasi dari perangkat lokal yang digunakan untuk implementasi:
\begin{itemize}
	\item Sistem Operasi: Ubuntu 24.04.4, berjalan di lingkungan WSL2
	\item CPU: Intel Core i5-11400H (3 core fisik dan 6 \textit{logical thread} dialokasikan untuk WSL2)
	\item Memori RAM: 16 GB (8 GB teralokasi untuk WSL2)
	\item \textit{Tools}:
	      \begin{itemize}
		      \item Rust (\texttt{rustc 1.93.0}, \texttt{cargo 1.93.0})
		      \item Node.js \texttt{v24.13.0}
		      \item Docker \texttt{29.5.3}
	      \end{itemize}
\end{itemize}

Sementara itu, spesifikasi kedua VPS yang digunakan adalah sebagai berikut:
\begin{itemize}
	\item Sistem Operasi: Ubuntu 24.04.4 LTS
	\item Jumlah vCPU: 8 vCPU
	\item Memori RAM: 8 GB
	\item \textit{Tools}:
	      \begin{itemize}
		      \item Docker: Docker \texttt{29.4.2}
	      \end{itemize}
\end{itemize}